\FloatBarrier
\section{Downstream shielding plate and vacuum jacket}
\label{sec:magnets-plate}

Upcoming in GEMC2 6.2 and CLAS12 GEMC Systems 0.4, DownstreamShieldingPlate and DownstreamVacuumJacket are
optimized as two separate volumes under their original names. Both retain their stainless steel, passive
180-degree y rotation, and 1273.27 mm z translation in the default, rga\_spring2018, and rga\_fall2018
configurations (runs 11, 3029, and 4763). All source and runtime copies match in GEMC2 and GEMC3; original
engineering inputs are archived in the owning GEMC2 repository.

The plate fills eight through-holes: six bolt holes of approximately 9.525 mm radius and two smaller holes of
approximately 2.38125 mm radius. The jacket fills six blind drilled recesses of approximately 6.35 mm entry
radius. Main circular rims and the hexagonal jacket step are extracted from the original vertices; only the
flat-face triangulation and fastener features change. No Poisson reconstruction, vertex displacement, or
joining of the separated parts is applied. The main aperture radii remain 76.35 mm for the plate and 63.5 mm
for the jacket, with the exact engineering polygon coordinates. The 32.83203125 mm axial gap and both original
bounds are retained.

The plate falls from 1,380 to 508 triangles (63.19\% fewer); the jacket falls from 2,392 to 556 (76.76\%
fewer). Together they fall from 3,772 to 1,064 triangles (71.79\% fewer) and from 188,768 to 53,368 bytes.
Each result is a single closed, consistently oriented manifold component of genus one, with no boundary edges,
detected self-intersections, or zero-area triangles. Independent polygon-area and thickness calculations match
the filled STL volumes within the stated builder tolerance of one part per billion. These checks describe the
stored surfaces, not a certified engineering mass.

Signed plate volume changes from 243534.863634 to 249067.081984 cubic millimetres (+2.271633\%); jacket volume
changes from 4082168.940709 to 4101145.339931 cubic millimetres (+0.464861\%). The pair gains 0.566581\%
signed steel volume from the authorized hole filling. This is a deliberate material-model change. The magnets
collection retains 51 meshes and now contains 129,522 triangles, after the separately documented
redundant-shielding removal.

Independent VTK containment checks identify all fourteen drilled features from the raw meshes. Forty-eight
probes that were inside the original holes are inside steel after filling; 1,296 bore probes remain outside
steel. Another 10,080 envelope probes away from the fastener features match the original geometry, and all 120
gap probes remain outside both separate volumes. These are sampled checks; exact rim-coordinate preservation
and analytic volume checks provide additional constraints.

The isolated benchmark loads only these two CAD rows at run 11, before and after optimization, with the same
fixed GEMC3 dev executable and Geant4 11.4.3 on macOS arm64. Both fixtures use FTFP\_BERT, one thread, default
production cuts, and no field, optical physics, GUI, or event output. Each 5 GeV/c primary starts at the
origin with polar angle uniform in solid angle from 0.5 to 3.2 degrees and full azimuth. There is no target,
luminosity normalization, or pileup. Five alternating-order repetitions use 2,000 events per particle and
seeds 12345--12349; LUND primary hashes match within every seed pair. All fifty timed processes and both
separate construction diagnostics complete, with all forty tracking processes reporting 2,000 events.

Table~\ref{tab:magnets-plate-tracking} separates tracking from the startup measurements.

Median zero-event process wall time is 0.73826 s before and 0.73375 s after; CPU time is 0.710 and 0.700 s,
and peak RSS is 124.05 and 120.06 MiB. Tracking values are median Geant4 beamOn wall times excluding startup.
Ratios are before/after: values below one indicate slower tracking. These observations describe the final
two-volume model; material changes and elimination of rejected source facets prevent interpreting the
comparison as a pure tessellation-speed measurement. It establishes neither full-CLAS12 performance nor
response equivalence.

\begin{table}[htbp]
\centering
\setlength{\tabcolsep}{8pt}
\caption{Separate downstream plate and jacket tracking; five repetitions of 2,000 events per particle.}
\label{tab:magnets-plate-tracking}
\begin{tabular}{lrrrr}
\toprule
Particle & Before (s) & After (s) & Speedup & Paired ratio range \\
\midrule
Electron & 14.25750 & 11.78500 & 1.210 & 1.199--1.223 \\
Proton & 5.53148 & 5.14433 & 1.075 & 1.061--1.107 \\
Positive pion & 5.51953 & 5.00997 & 1.102 & 1.089--1.103 \\
Neutron & 3.94107 & 3.68135 & 1.071 & 1.068--1.117 \\
\bottomrule
\end{tabular}
\end{table}

Each isolated raw import reports ten GeomSolids1001 and four GeomSolids1002 warning blocks; each optimized
import is clean. The optimized timed tracking processes and isolated check\_overlaps=2 construction report no
Geant4 exceptions. All three runs pass standalone GEMC3 magnets construction (53 volumes including the world),
combined GEMC2 TEXT/CAD and SQLite/CAD initialization, and native/CAD semantic checks with one solenoid row
and 51 CAD rows. Existing full-magnets assembly and unrelated mesh warnings remain in the retained validation
logs. The run-4763 GEMC3 beamline-bank limitation documented in the preceding section also remains. Isolated
construction success does not certify the complete assembly or detector response.

Figure~\ref{fig:magnets-plate} compares the original and optimized surfaces.
\begin{figure}[p]
\centering
\includegraphics[width=0.98\linewidth,height=0.76\textheight,keepaspectratio]
{figures/downstream_plate_comparison.pdf}
\caption{
Original (left) and optimized (right) stored STL surfaces in millimetres, with common bounds, orthographic
projection and physical scale within each row. The first row shows the plate entry face, the second the
drilled jacket face, and the third both separately placed parts at their original spacing. The face panels
show stored planar facets; assembly panels use complete stored meshes, including triangles rejected by Geant4.}
\label{fig:magnets-plate}
\end{figure}

The dedicated data/magnets\_plate directory retains the mesh inventory, independent feature checks, all timing
samples, paired-input provenance, and three-run validation. magnets\_plate\_archive.zip contains both raw and
optimized two-volume fixtures, raw logs and process reports, source snapshots, matched primary files, original
engineering meshes, and reproduction scripts. The fixed executable is identified by hash and is not
redistributed. Earlier study archives remain unchanged.
